\section{Señal de entrada y coeficientes}\label{sec:senales}

\subsection{Símbolos QPSK y sobremuestreo}

Cada símbolo toma dos bits $(b_I, b_Q)$, con bit 0 $\to +1$ y bit 1 $\to -1$, sin normalizar por $\sqrt{2}$: $a[m] \in \{\pm1 \pm j\}$. Un sobremuestreador 2×, separado del generador, intercala un cero después de cada símbolo:
\begin{equation}
  x[2m] = a[m], \qquad x[2m+1] = 0,
\end{equation}
de modo que $S$ símbolos producen $N = 2S$ muestras. Ambas ramas reciben exactamente las mismas muestras. La tabla~\ref{tab:qpsk} muestra el comienzo de la secuencia de la semilla 1.

\begin{table}[h]
  \centering
  \caption{Secuencia QPSK de ejemplo (semilla 1).}\label{tab:qpsk}
  \begin{tabular}{rcccc}
    \toprule
    $m$ & $b_I\,b_Q$ & $a[m]$ & $x[2m]$ & $x[2m+1]$ \\
    \midrule
    0 & 0\,1 & $+1-j$ & $+1-j$ & 0 \\
    1 & 1\,1 & $-1-j$ & $-1-j$ & 0 \\
    2 & 0\,0 & $+1+j$ & $+1+j$ & 0 \\
    3 & 1\,1 & $-1-j$ & $-1-j$ & 0 \\
    \bottomrule
  \end{tabular}
\end{table}

\subsection{Coeficientes RRCOS}

El roll-off y el sobremuestreo no determinan una tabla única: falta elegir la fase de muestreo y la normalización. Se muestreó el pulso RRCOS en forma cerrada ($\alpha = 0{,}5$, $t$ en unidades del período de símbolo $T_s$)
\begin{equation}
  p(t) = \frac{\sin\!\big(\pi t(1-\alpha)\big) + 4\alpha t \cos\!\big(\pi t(1+\alpha)\big)}
              {\pi t\,\big(1 - (4\alpha t)^2\big)}
\end{equation}
en la grilla simétrica $t_k = (k - 3{,}5)/2$, $k = 0\ldots7$, sin coeficiente central, y se normalizó a energía unitaria:
\begin{equation}
  h[k] = \frac{p(t_k)}{\sqrt{\sum_j p(t_j)^2}}, \qquad
  y[n] = \sum_{k=0}^{7} h[k]\, x[n-k].
\end{equation}
Los límites $t = 0$ y $t = \pm 1/(4\alpha)$ se evalúan con sus expresiones cerradas. La tabla~\ref{tab:coef} da los valores. El filtro es simétrico, de fase lineal y con retardo de grupo de 3,5 muestras; su ganancia en continua es $\sum h = 1{,}4007$ y la salida cumple $|y| \le \sum|h| = 1{,}8445$ por componente.

\begin{table}[h]
  \centering
  \caption{Coeficientes del filtro y su versión cuantizada en Q1.7 (sección~\ref{sec:punto-fijo}).}\label{tab:coef}
  \begin{tabular}{rrrr}
    \toprule
    $k$ & $t_k/T_s$ & $h[k]$ & $h_q[k]$ \\
    \midrule
    0 & $-1{,}75$ & $+0{,}010942691568693$ & 1 \\
    1 & $-1{,}25$ & $-0{,}110955576454819$ & $-14$ \\
    2 & $-0{,}75$ & $+0{,}110955576454819$ & 14 \\
    3 & $-0{,}25$ & $+0{,}689389568827662$ & 88 \\
    4 & $+0{,}25$ & $+0{,}689389568827662$ & 88 \\
    5 & $+0{,}75$ & $+0{,}110955576454819$ & 14 \\
    6 & $+1{,}25$ & $-0{,}110955576454819$ & $-14$ \\
    7 & $+1{,}75$ & $+0{,}010942691568693$ & 1 \\
    \bottomrule
  \end{tabular}
\end{table}

Con 8 coeficientes a 2× el pulso abarca solo cuatro símbolos, por lo que la atenuación fuera de banda es pobre. Es consecuencia de combinar los 8 coeficientes que fija la consigna con el sobremuestreo 2× elegido por el equipo, no de la cuantización.
